\chapter{What the Network Hands You, and Where Those Numbers Physically Come From}
\label{ch:network_outputs}

You show the board to the neural network \textbf{once}. One forward pass. It returns \textbf{two} things, and it is important to understand that they are two different heads of the same network answering two different questions.

\section{The Value Head: Position Outcome Probabilities}

The value head answers: \textit{``How is this position going for the side to move?''}

It outputs three probabilities:
\begin{align*}
  w &= \text{probability this position ends in a win for the side to move} \\
  d &= \text{probability this position ends in a draw} \\
  l &= \text{probability this position ends in a loss for the side to move} \\[1ex]
  w + d + l &= 1 \quad \text{(a game must end in exactly one of the three)}
\end{align*}

\begin{established}[Where These Numbers Physically Come From]
They are not dice. Chess is deterministic --- the true theoretical outcome of a position is exactly one of win, draw, or loss. These probabilities are the network's \textit{degree of belief}, learned from millions of self-play games. Read $w = 0.60$ as:
\begin{quote}
``Among the positions I was trained on that looked like this one, the side to move went on to win about 60\% of the time.''
\end{quote}
That is \textbf{epistemic} uncertainty --- uncertainty from limited knowledge, not from randomness in the game.
\end{established}

Two derived numbers:
\begin{align*}
  \text{Expected score} \quad E &= w + 0.5 \cdot d \quad \text{range } [0.0, 1.0] \\
  \text{Net value} \quad V &= w - l \quad \text{range } [-1.0, +1.0]
\end{align*}
and they are the exact same information rescaled:
\begin{equation}
  E = 0.5 + 0.5 \cdot V \iff V = 2 \cdot E - 1
\end{equation}

For our position the network returns \textbf{$V = +0.97602$}.

\begin{center}
% fig_c_wdl.tex -- WDL simplex and aligned E/V scales
\begin{tikzpicture}[scale=0.95, every node/.style={font=\small}]
  % Stacked WDL bar
  \pgfmathsetmacro{\barW}{10.0}
  \pgfmathsetmacro{\wW}{6.0}   % w = 0.60
  \pgfmathsetmacro{\dW}{3.0}   % d = 0.30
  \pgfmathsetmacro{\lW}{1.0}   % l = 0.10

  \fill[IdeaGreen!70] (0, 0) rectangle (\wW, 0.9);
  \fill[WorkGrey!35] (\wW, 0) rectangle (\wW+\dW, 0.9);
  \fill[PitfallRed!65] (\wW+\dW, 0) rectangle (\barW, 0.9);
  \draw[NavyBlue, line width=0.8pt] (0,0) rectangle (\barW, 0.9);
  \draw[NavyBlue, line width=0.6pt] (\wW, 0) -- (\wW, 0.9);
  \draw[NavyBlue, line width=0.6pt] (\wW+\dW, 0) -- (\wW+\dW, 0.9);

  % Labels on bar segments
  \node[white, font=\footnotesize\bfseries] at (\wW/2, 0.45) {$w = 0.60$ (Win)};
  \node[black, font=\footnotesize\bfseries] at (\wW + \dW/2, 0.45) {$d = 0.30$ (Draw)};
  \node[white, font=\tiny\bfseries] at (\wW + \dW + \lW/2, 0.45) {$l{=}0.10$};

  % E marker at w + 0.5*d = 6.0 + 1.5 = 7.5cm
  \draw[NavyBlue, line width=1.2pt, -{Stealth[length=2.5mm]}] (7.5, 1.6) -- (7.5, 0.95);
  \node[above, font=\footnotesize\bfseries, text=NavyBlue] at (7.5, 1.6) {Expected Score: $E = w + \frac{1}{2}d = 0.75$};
  \draw[NavyBlue, dashed, line width=0.6pt] (7.5, 0.9) -- (7.5, -2.4);

  % Axis 1: Expected Score E in [0, 1]
  \draw[WorkGrey, line width=0.7pt, ->] (0, -0.8) -- (10.5, -0.8) node[right, font=\footnotesize] {$E \in [0, 1]$};
  \foreach \x/\val in {0/0.0, 5/0.5, 7.5/0.75, 10/1.0} {
    \draw[WorkGrey, line width=0.6pt] (\x, -0.7) -- (\x, -0.9);
    \node[below, font=\footnotesize] at (\x, -0.9) {\val};
  }

  % Axis 2: Net Value V in [-1, +1]
  \draw[WorkGrey, line width=0.7pt, ->] (0, -2.0) -- (10.5, -2.0) node[right, font=\footnotesize] {$V \in [-1, +1]$};
  \foreach \x/\val in {0/-1.0, 5/0.0, 7.5/+0.50, 10/+1.0} {
    \draw[WorkGrey, line width=0.6pt] (\x, -1.9) -- (\x, -2.1);
    \node[below, font=\footnotesize\bfseries, text=NavyBlue] at (\x, -2.1) {\val};
  }

  % Formula bridge
  \node[anchor=west, font=\footnotesize\itshape, text=WorkGrey] at (0, -2.9) {%
    Relationship: $V = w - l = 0.60 - 0.10 = \mathbf{+0.50} \iff V = 2E - 1 = 2(0.75) - 1 = \mathbf{+0.50}$%
  };
\end{tikzpicture}

\captionof{figure}{\textbf{The WDL outcome simplex and aligned $E$ and $V$ scales.} Expected score $E = w + 0.5d$ and Net value $V = w - l$ represent the same information centered at $0.5$ and $0.0$ respectively.}
\label{fig:wdl}
\end{center}

\begin{center}
% fig_c_sharp.tex -- Sharp vs Drawish comparison at identical E = 0.60
\begin{tikzpicture}[scale=0.85, every node/.style={font=\small}]
  % Position A: Sharp Sicilian
  \node[font=\footnotesize\bfseries, text=NavyBlue, anchor=south] at (2.2, 3.6) {Position A: Sharp Tactical Brawl};
  \fill[IdeaGreen!70] (0, 0) rectangle (4.4, 0.55*3.5);
  \fill[WorkGrey!35] (0, 0.55*3.5) rectangle (4.4, 0.65*3.5);
  \fill[PitfallRed!65] (0, 0.65*3.5) rectangle (4.4, 1.0*3.5);
  \draw[NavyBlue, line width=0.8pt] (0,0) rectangle (4.4, 3.5);
  \draw[NavyBlue, line width=0.5pt] (0, 0.55*3.5) -- (4.4, 0.55*3.5);
  \draw[NavyBlue, line width=0.5pt] (0, 0.65*3.5) -- (4.4, 0.65*3.5);

  \node[white, font=\tiny\bfseries] at (2.2, 0.95) {$w = 0.55$ (55\% Win)};
  \node[black, font=\tiny\bfseries] at (2.2, 2.1) {$d = 0.10$ (10\% Draw)};
  \node[white, font=\tiny\bfseries] at (2.2, 2.9) {$l = 0.35$ (35\% Loss)};

  \node[font=\footnotesize, align=center, anchor=north] at (2.2, -0.2) {%
    $E = 0.55 + \frac{1}{2}(0.10) = \mathbf{0.60}$\\[0.3ex]
    \textit{Razor sharp: high volatility}%
  };

  % Position B: Dry Rook Endgame
  \begin{scope}[xshift=6cm]
    \node[font=\footnotesize\bfseries, text=NavyBlue, anchor=south] at (2.2, 3.6) {Position B: Dry Rook Endgame};
    \fill[IdeaGreen!70] (0, 0) rectangle (4.4, 0.20*3.5);
    \fill[WorkGrey!35] (0, 0.20*3.5) rectangle (4.4, 1.0*3.5);
    \draw[NavyBlue, line width=0.8pt] (0,0) rectangle (4.4, 3.5);
    \draw[NavyBlue, line width=0.5pt] (0, 0.20*3.5) -- (4.4, 0.20*3.5);

    \node[white, font=\tiny\bfseries] at (2.2, 0.35) {$w = 0.20$ (20\% Win)};
    \node[black, font=\tiny\bfseries] at (2.2, 2.1) {$d = 0.80$ (80\% Draw)};

    \node[font=\footnotesize, align=center, anchor=north] at (2.2, -0.2) {%
      $E = 0.20 + \frac{1}{2}(0.80) = \mathbf{0.60}$\\[0.3ex]
      \textit{Completely safe: $l = 0.00$}%
    };
  \end{scope}

  % Summary note below
  \node[draw=NavyBlue!40, fill=PaleGrey, rounded corners=3pt, inner sep=4pt, font=\scriptsize, align=center] at (5.2, -1.8) {%
    \textbf{The Key Insight:} Both positions have identical Expected Score ($E = 0.60$, $\approx +0.80$ cp),\\
    yet their character is completely different. The full $(w, d, l)$ triad captures this instantly.%
  };
\end{tikzpicture}

\captionof{figure}{\textbf{Why WDL matters: sharp vs.\ drawish positions.} Two positions with identical expected score $E = 0.60$ have completely different volatility and risk profiles.}
\label{fig:sharp}
\end{center}

\paragraph{Why the search uses $V$ and not $E$.}
Chess is zero-sum. Whatever is good for White by exactly that much is bad for Black. On the $V$ scale, switching sides is a single multiplication:
\begin{equation}
  V_{\text{from the other side}} = -V
\end{equation}
On the $E$ scale you would have to write $1 - E$, which is clumsier to propagate. That sign flip is the whole reason $V$ exists.

\section{The Policy Head: Move Attention Distribution}

The policy head answers: \textit{``Which moves look worth examining?''}

For our position:
\begin{center}
\begin{tabular}{lrr}
  \toprule
  Move & Policy Prior $P$ & Percentage \\
  \midrule
  \textbf{Kd6} & 0.4513 & 45.13\% \\
  \textbf{Kf6} & 0.4423 & 44.23\% \\
  \textbf{Kf5} & 0.0538 & 5.38\% \\
  \textbf{Kd5} & 0.0526 & 5.26\% \\
  \midrule
  \textbf{Total} & & 100.00\% \\
  \bottomrule
\end{tabular}
\end{center}

\begin{established}[What $P$ Is Not]
$P$ is \textit{not} the evaluation of the move. The network has not looked at any of these positions. $P$ is a trained reflex --- ``in positions of this shape, these are the moves that turned out to be worth examining.'' It is attention, not judgement.
\end{established}

Notice the network put 89.4\% of its attention on the two winning king moves and still left 10.6\% on the two that throw the win away. It is not certain. Good --- that residue is what lets the search check them cheaply and move on.
